\chapter{旋量、自旋与双群}
\label{chap:spinors-double-groups}

$SO(3)$ 的普通矢量在 $2\pi$ 转动后回到自身，半整数自旋态却获得负号。上一章已经用 $SU(2)$ 的轴角指数看见了这个差别；现在需要给承载这种变换的二分量复向量一个明确身份。

旋量不是空间矢量的另一种坐标写法，而是 $SU(2)$ 的基本表示。把它限制到分子的离散转动对称时，$2\pi$ 转动仍不能与恒等操作混同，因此普通点群必须提升为双群。自旋--轨道耦合的表示分类由此获得一致的群论语言。

\section{旋量、双群与自旋--轨道耦合}
\label{sec:double-groups-spinors}


对整数 $j$，$2\pi$ 转动表示为 $+I$，所以态可以直接承载 $SO(3)$ 或普通点群的表示；对半整数 $j$，却有
\[
U(2\pi)=-I.
\]
这并不表示物理态在 $2\pi$ 后变成“另一种可观测状态”：量子态整体相位本身不可观测。但如果体系中需要比较不同路径、不同对称操作的相位，或者要构造群表示、分析自旋--轨道耦合，负号就不能被随意丢弃。数学上最干净的做法是把每个空间转动提升到 $SU(2)$ 中的两个元素，并保留中央元素
\[
\bar E=-I_2,
\qquad \bar E^2=E.
\]
它代表“空间上转了 $2\pi$，但在旋量空间中并非恒等”的操作。

\begin{definition}[旋转点群的双群]
设 $G\subset SO(3)$ 为有限纯转动点群，$\pi:SU(2)\to SO(3)$ 为二重覆盖。定义
\[
G^D=\pi^{-1}(G)\subset SU(2).
\]
则 $G^D$ 称为 $G$ 的双群。它的阶为
\[
|G^D|=2|G|,
\]
并含中央元素 $\bar E=-I_2$。普通表示满足 $D(\bar E)=+I$，旋量表示满足 $D(\bar E)=-I$。
\label{def:double-group}
\end{definition}

对包含反演或反射的点群，实际双群要在上述旋转子群的提升之外再加入相应非纯转动操作的提升。应用时最重要的不是抽象构造细节，而是两类不可约表示的区别：单值表示看不见 $\bar E$，双值表示在 $\bar E$ 下取负号。附录中给出常用双群的记号与特征标使用规则。

电子自旋 $1/2$ 的标准旋量
\[
\chi=
\begin{pmatrix}c_+\\c_-\end{pmatrix}
\in\CC^2
\]
在转动下按 $j=\tfrac12$ 表示变换，
\[
\chi\longmapsto
U(\uv n,\theta)\chi
=
\left(
\cos\frac\theta2 I_2
-\ii\sin\frac\theta2\,\uv n\cdot\boldsymbol{\sigma}
\right)\chi.
\]
特别地，绕 $z$ 轴转动角 $\theta$ 时
\[
\ket{\uparrow}\mapsto \ee^{-\ii\theta/2}\ket{\uparrow},
\qquad
\ket{\downarrow}\mapsto \ee^{\ii\theta/2}\ket{\downarrow}.
\]
这就是旋量“半角”行为的最直接表现。

在中心势近似下，轨道角动量 $\vb L$ 与自旋 $\vb S$ 可以先分别讨论。相对论修正产生的自旋--轨道项具有一般形式
\begin{equation}
H_{\mathrm{SO}}=\lambda(r)\,\vb L\cdot\vb S,
\label{eq:spin-orbit-hamiltonian}
\end{equation}
或者在给定壳层中用有效常数 $\lambda$ 近似。定义总角动量
\[
\vb J=\vb L+\vb S.
\]
由
\[
J^2=L^2+S^2+2\vb L\cdot\vb S
\]
得到
\begin{equation}
\vb L\cdot\vb S
=\frac12\bigl(J^2-L^2-S^2\bigr).
\label{eq:LdotS}
\end{equation}
因此在 $|\ell,s;j,m\rangle$ 基中，自旋--轨道耦合已经对角化：
\begin{equation}
E_{\mathrm{SO}}(j)
=\frac{\lambda\hbar^2}{2}
\bigl[j(j+1)-\ell(\ell+1)-s(s+1)\bigr].
\label{eq:spin-orbit-energy}
\end{equation}
对单电子 $s=\tfrac12$，允许
\[
j=\ell\pm\frac12.
\]
于是同一个轨道壳层会分成两个总角动量多重态。这一结论常在原子物理里直接作为角动量耦合规则出现；从群论角度看，它来自
\[
D^{(\ell)}\otimes D^{(1/2)}
\cong D^{(\ell+1/2)}\oplus D^{(\ell-1/2)}.
\]
下一节将把这个 Clebsch--Gordan 分解从一般表示理论推出。

在晶体中，完整 $SO(3)$ 已经先被晶场降低到有限点群 $G$。若暂时忽略自旋--轨道耦合，轨道态按普通点群不可约表示 $\Gamma_{\mathrm{orb}}$ 分类，而自旋还保留独立的二维空间。打开自旋--轨道耦合后，真正要约化的是
\begin{equation}
\Gamma_{\mathrm{orb}}\otimes\Gamma_{1/2}
\end{equation}
作为双群 $G^D$ 的表示。于是原先“每个轨道态再乘一个自旋二重简并”的图像会变成双群不可约表示的分裂与重组。

\begin{example}[$O_h$ 晶场中的 $d$ 轨道与自旋]
无自旋时，$\ell=2$ 的五维 $SO(3)$ 表示限制到 $O_h$ 后分解为
\[
D^{(2)}\downarrow_{O_h}
\cong E_g\oplus T_{2g},
\]
维数为 $2+3=5$。加入自旋 $1/2$ 后，总空间维数变为 $10$。在常用 Koster 记号中，基本旋量表示记为 $\Gamma_6$，有
\[
E_g\otimes\Gamma_6\cong\Gamma_8,
\qquad
T_{2g}\otimes\Gamma_6\cong\Gamma_7\oplus\Gamma_8,
\]
维数分别为 $2\times2=4$ 与 $3\times2=2+4$。因此双群并不是简单地把原来的 $E_g,T_{2g}$ 各复制两份，而是重新组织成旋量不可约表示。具体能量顺序仍取决于晶场和自旋--轨道耦合常数；群论决定的是允许怎样分裂，而不是自动给出数值能量。
\end{example}

进一步分析双群特征标表的系统构造。

上面的例子用了 Koster 记号 $\Gamma_6,\Gamma_7,\Gamma_8$，但没有说明这些表示和特征标从何而来。现在给出系统方法。

\begin{proposition}[双群特征标表的构造方法]
\label{prop:double-group-char}
设 $G$ 是晶体点群，$G^D=G\cup\bar EG$ 是其双群。$G^D$ 的不可约表示分两类：
\begin{enumerate}[label=(\roman*),leftmargin=3em]
\item \textbf{单值表示}（$D(\bar E)=+I$）：与 $G$ 的不可约表示一一对应，特征标 $\chi^{D^+}(R)=\chi^G(R)$，$\chi^{D^+}(\bar ER)=\chi^G(R)$。
\item \textbf{双值表示}（$D(\bar E)=-I$）：额外的新表示，特征标 $\chi^{D^-}(R)=-\chi^{D^-}(\bar ER)$。双值表示只在 $G^D$ 中出现，不下降到 $G$。
\end{enumerate}
双值表示的个数和维数由维数平方和公式确定：若 $G$ 有 $r$ 个不可约表示（维数 $d_1,\dots,d_r$），$G^D$ 另有 $s$ 个双值不可约表示（维数 $\tilde d_1,\dots,\tilde d_s$），则
\[
|G^D|=2|G|=\sum_{i=1}^r d_i^2+\sum_{j=1}^s \tilde d_j^2,
\]
其中 $\tilde d_j$ 是双值表示的维数。单值表示贡献 $\sum_{i=1}^r d_i^2=|G|$，故 $\sum_{j=1}^s\tilde d_j^2=|G|$。
\end{proposition}

\begin{proof}
单值表示。$G^D/G\cong\{\bar E^0,\bar E^1\}\cong C_2$（因 $\bar E$ 是 $G^D$ 的中心正规子群）。$G$ 的每个不可约表示 $\rho$ 提升为 $G^D$ 的单值表示：$\tilde\rho(R)=\rho(R)$，$\tilde\rho(\bar ER)=\rho(R)$（即 $\tilde\rho(\bar E)=I$）。这给出 $r$ 个单值表示。

双值表示。$G^D$ 的其他表示满足 $\tilde\rho(\bar E)=-I$（因 $\bar E^2=E$ 蕴含 $\tilde\rho(\bar E)^2=I$，故 $\tilde\rho(\bar E)=\pm I$；$\tilde\rho(\bar E)=+I$ 已是单值表示）。$\tilde\rho(\bar ER)=-\tilde\rho(R)$，故 $\chi(\bar ER)=-\chi(R)$。

维数公式。$|G^D|=2|G|=\sum_{\text{所有不可约}} d_i^2$。单值表示贡献 $\sum_{i=1}^r d_i^2=|G|$，故双值表示贡献 $\sum\tilde d_j^2=|G|$。双值表示成对出现（若 $\tilde\rho$ 是双值表示，$\tilde\rho\otimes\operatorname{sgn}$（$\operatorname{sgn}$ 是 $G^D/G\cong C_2$ 的符号表示）也是双值表示），但通常只计独立的一个。
\end{proof}

\medskip
\noindent\emph{$O^D$ 的双群特征标表。}
正八面体群 $O$（24 元素，5 类：$E,8C_3,3C_2,6C_4,6C_2'$）的双群 $O^D$ 有 48 元素。$O^D$ 的共轭类：$E$ 与 $\bar E$ 各自成一类；$8C_3$ 与 $8\bar EC_3$ 分裂；$6C_4$ 与 $6\bar EC_4$ 分裂。面轴 $\pi$ 转动 $C_2$ 与 $\bar EC_2$ 共轭、合并为一类（6 个元素），边轴 $\pi$ 转动 $C_2'$ 与 $\bar EC_2'$ 也合并为一类（12 个元素）。最终 $O^D$ 有 8 个共轭类：
\[
E,\;\bar E,\;8C_3,\;8\bar EC_3,\;6C_4,\;6\bar EC_4,\;6C_2,\;12C_2'.
\]
各类大小相加 $1+1+8+8+6+6+6+12=48=|O^D|$。

由 $|O|=24$，双值表示满足 $\sum\tilde d_j^2=24$。$O$ 的单值表示维数 $1,1,2,3,3$（$A_1,A_2,E,T_1,T_2$），贡献 $1+1+4+9+9=24$。双值表示维数满足 $\tilde d_1^2+\tilde d_2^2+\tilde d_3^2=24$，唯一解 $\tilde d=2,2,4$（因 $4+4+16=24$）。三个双值表示 $\Gamma_6,\Gamma_7,\Gamma_8$（Koster 记号），维数 $2,2,4$。

双值表示的特征标利用 $\chi(\bar ER)=-\chi(R)$ 和正交性确定。对 $j=1/2$（基本旋量），$\chi_{\Gamma_6}(R)=\tr D^{(1/2)}(R)$：$E\to2$，$C_3\to1$，$C_2\to0$，$C_4\to\sqrt2$，$C_2'\to0$。对 $j=3/2$，$\chi_{\Gamma_8}(R)=\tr D^{(3/2)}(R)$：$E\to4$，$C_3\to-1$，$C_2\to0$，$C_4\to0$，$C_2'\to0$。$\Gamma_7$ 由正交性确定。

完整的 $O^D$ 双值表示特征标表（仅列 $G$ 部分，$\bar EG$ 部分取负）：
\begin{center}
\begin{tabular}{c|ccccc}
\toprule
$O^D$ & $E$ & $8C_3$ & $3C_2$ & $6C_4$ & $6C_2'$ \\
\midrule
$\Gamma_6$ & $2$ & $1$ & $0$ & $\sqrt2$ & $0$ \\
$\Gamma_7$ & $2$ & $1$ & $0$ & $-\sqrt2$ & $0$ \\
$\Gamma_8$ & $4$ & $-1$ & $0$ & $0$ & $0$ \\
\bottomrule
\end{tabular}
\end{center}
验证正交性：$\langle\Gamma_6,\Gamma_6\rangle=\frac1{24}(4+8+0+12+0)=1$ $\checkmark$。$\langle\Gamma_6,\Gamma_7\rangle=\frac1{24}(4+8+0-12+0)=0$ $\checkmark$。$\langle\Gamma_8,\Gamma_8\rangle=\frac1{24}(16+8+0+0+0)=1$ $\checkmark$。
\par\medskip

\medskip
\noindent\emph{用双群特征标预测晶场劈裂。}
$d$ 轨道（$\ell=2$）在 $O$ 晶场中劈裂为 $E\oplus T_2$。加入自旋 $1/2$ 后，总表示为 $(E\oplus T_2)\otimes\Gamma_6$。用双群特征标计算：

$E\otimes\Gamma_6$ 的特征标 $\chi_E\cdot\chi_{\Gamma_6}$（$\chi_E=(-1)$ 在 $C_3$ 上）为 $(4,-1,0,0,0)$（在 $E,8C_3,3C_2,6C_4,6C_2'$ 上）。分解：
\[
\langle\chi_E\chi_{\Gamma_6},\Gamma_8\rangle=\frac1{24}(4\cdot4+8\cdot(-1)\cdot(-1)+0+0+0)=\frac{24}{24}=1,
\]
\[
\langle\chi_E\chi_{\Gamma_6},\Gamma_6\rangle=\frac1{24}(4\cdot2+8\cdot(-1)\cdot1+0+0+0)=\frac{0}{24}=0.
\]
故 $E\otimes\Gamma_6\cong\Gamma_8$（四维）。同理 $T_2\otimes\Gamma_6$ 的特征标为 $(6,0,0,-\sqrt2,0)$（$T_2$ 在 $C_4$ 上为 $-1$），分解为 $\Gamma_7\oplus\Gamma_8$（$2+4=6$ 维）。总分解 $\Gamma_8\oplus\Gamma_7\oplus\Gamma_8$，维数 $4+2+4=10=5\times2$ $\checkmark$。

Kramers 简并验证。双值表示 $\Gamma_6$（维 2）、$\Gamma_7$（维 2）、$\Gamma_8$（维 4）的维数都是偶数——这正是 Kramers 定理的体现：半整数自旋系统的每个能级至少二重简并。$\Gamma_8$ 是四重简并（两个 Kramers 对），$\Gamma_6$ 和 $\Gamma_7$ 各是一个 Kramers 对。在 $O_h$ 中加反演 $g/u$ 下标后，$d$ 轨道（偶宇称）的旋量表示为 $\Gamma_8^+\oplus\Gamma_7^+\oplus\Gamma_8^+$。
\par\medskip

晶体场中自旋--轨道耦合的分支图如图 \figref{fig:crystal-spin-orbit-branching} 所示。
\begin{figure}[htbp]
\centering
\begin{tikzfig}
  \node[lecturebox] (so3) at (0,2.0) {$SO(3)$\\$d:\ \ell=2$};
  \node[lecturebox] (oh) at (4.3,2.0) {$O_h$\\$E_g\oplus T_{2g}$};
  \node[lecturebox] (double) at (8.6,2.0) {$O_h^D$\\$\Gamma_8\oplus\Gamma_7\oplus\Gamma_8$};
  \draw[lecturearrow] (so3)--node[above,lecturelabel]{晶场降对称}(oh);
  \draw[lecturearrow] (oh)--node[above,lecturelabel]{乘 $\Gamma_{1/2}$ 并约化}(double);
  \node[lecturelabel,align=center] at (4.3,.25)
  {普通点群首先整理轨道简并；\\双群再把自旋与轨道作为一个整体分类。};
\end{tikzfig}
\caption{晶场与自旋--轨道耦合对应的两级表示约化。}
\label{fig:crystal-spin-orbit-branching}
\end{figure}
